\section{Laplace Transform}

\subsection{Definitions}

\subsubsection{Unilateral Laplace transform}
\[
X(s)=\mathcal{L}\{x(t)\}=\int_{0^-}^{\infty}x(t)e^{-st}\,dt.
\]
The exams state that Laplace transforms are unilateral unless otherwise specified. Initial conditions are included through derivative rules.

\subsubsection{Inverse Laplace transform}
\[
x(t)=\mathcal{L}^{-1}\{X(s)\}.
\]
Inverse transformation is normally performed by partial fractions and transform pairs.

\subsubsection{Region of convergence}
\[
X(s)=\int x(t)e^{-st}\,dt \quad \text{converges for } s \in \mathrm{ROC}.
\]
For causal rational signals, the region of convergence lies to the right of the rightmost pole.

\subsection{Transform Pairs}

\subsubsection{Step, impulse, and ramp Laplace transforms}
\[
\delta(t)\leftrightarrow 1,\qquad
u(t)\leftrightarrow \frac{1}{s},\qquad
t u(t)\leftrightarrow \frac{1}{s^2}.
\]
These pairs are the basic inputs for impulse, step, and ramp responses.

\subsubsection{Exponential Laplace transform}
\[
e^{-at}u(t)\leftrightarrow \frac{1}{s+a}.
\]
First-order poles produce causal exponentials in time.

\subsubsection{Damped sinusoid Laplace transforms}
\[
e^{-\alpha t}\cos\omega_dt\,u(t)\leftrightarrow
\frac{s+\alpha}{(s+\alpha)^2+\omega_d^2},
\]
\[
e^{-\alpha t}\sin\omega_dt\,u(t)\leftrightarrow
\frac{\omega_d}{(s+\alpha)^2+\omega_d^2}.
\]
Complex conjugate poles produce exponentially damped sinusoids.

\subsection{Laplace Properties}

\subsubsection{Linearity of Laplace transform}
\[
\mathcal{L}\{a x(t)+b y(t)\}=aX(s)+bY(s).
\]
Linearity allows transform expressions to be assembled term by term.

\subsubsection{Laplace derivative rule}
\[
\mathcal{L}\left\{\frac{d x}{dt}\right\}=sX(s)-x(0^-),
\]
\[
\mathcal{L}\left\{\frac{d^2 x}{dt^2}\right\}=s^2X(s)-s x(0^-)-\dot{x}(0^-).
\]
Unilateral derivative rules insert initial conditions directly into algebraic equations.

\subsubsection{Laplace integration rule}
\[
\mathcal{L}\left\{\int_{0^-}^{t}x(\tau)\,d\tau\right\}=\frac{X(s)}{s}.
\]
Time integration corresponds to division by \(s\).

\subsubsection{Laplace time shift}
\[
x(t-t_0)u(t-t_0)\leftrightarrow e^{-s t_0}X(s),\qquad t_0>0.
\]
A causal delay multiplies the Laplace transform by an exponential factor.

\subsubsection{Laplace convolution property}
\[
x(t)*h(t)\leftrightarrow X(s)H(s).
\]
Convolution in time is multiplication in the Laplace domain.

\subsubsection{Initial value theorem}
\[
x(0^+)=\lim_{s\to\infty}sX(s).
\]
The theorem applies when the limit exists and no impulse at the origin invalidates the ordinary value.

\subsubsection{Final value theorem}
\[
\lim_{t\to\infty}x(t)=\lim_{s\to 0}sX(s).
\]
The theorem applies when all poles of \(sX(s)\) are in the open left half-plane, except possibly a simple pole at the origin in \(X(s)\).
